\chapter{Model Categories and Projective Complexes}
\label{ch:model-categories}

A model category organizes localization by giving two compatible
procedures for changing an object or a map. We develop the homotopy
category and the Quillen adjunction directly from lifting, then establish
the unbounded projective model on complexes. Throughout, maps of
complexes have degree zero, and cochain differentials raise degree.
The conventions for model categories are those of \cite[Chapter 1]{hovey}
and \cite[Section 11.3]{riehl}.

\section{Lifting and homotopy}

For two maps \(i:A\to B\) and \(p:X\to Y\), write \(i\perp p\)
when every commutative square with \(i\) on the left and \(p\) on
the right has a diagonal filler \(B\to X\). A retract of arrows is
a retract in the category of commutative squares.

\begin{definition}[Model structure]\label{res:Q01}
A model category in this chapter is a complete and cocomplete category
with classes \(W\), \(\operatorname{Cof}\), and \(\operatorname{Fib}\)
of weak equivalences, cofibrations, and fibrations. The classes contain
isomorphisms and are closed under retracts. The class \(W\) satisfies
two-out-of-three. Call a map \emph{trivial} when it belongs to \(W\).
Cofibrations lift against trivial fibrations, and trivial cofibrations
lift against fibrations. Each map has functorial factorizations as a
cofibration followed by a trivial fibration and as a trivial cofibration
followed by a fibration.
Write \(0\) and \(1\) for the initial and terminal objects. An object
\(A\) is cofibrant if \(0\to A\) is a cofibration and fibrant if
\(A\to1\) is a fibration.
\end{definition}

The two lifting pairs determine their classes by orthogonality. Indeed,
if \(h\perp\operatorname{Fib}\), factor \(h=p i\) with \(i\) a
trivial cofibration and \(p\) a fibration. The lift against \(p\)
exhibits \(h\) as a retract of \(i\). The argument for the other
left class uses the other factorization, and duality gives the
right classes. A left class so determined is closed under pushout,
composition, transfinite composition, and retract. For a pushout,
the lift for the old map extends by the defining pushout. For an
ordinal sequence choose lifts successively and use the defining
colimit at limits. A right class has the dual closures, in
particular pullback and composition. These arguments supply the
closure properties used below.

\begin{proposition}[Cylinder and path objects]\label{res:Q02}
A cylinder for \(A\) is a factorization
\[
 A\amalg A\xrightarrow{(j_0,j_1)} C\xrightarrow{r}A
\]
of the fold map with first map a cofibration and \(r\in W\).
It is called good when \(r\) is a trivial fibration.
A path object for \(B\) is a factorization
\[
 B\xrightarrow{s} P\xrightarrow{(e_0,e_1)}B\times B
\]
of the diagonal with \(s\in W\) and second map a fibration.
It is called good when \(s\) is a trivial cofibration.
Good choices exist functorially. If \(A\) is cofibrant, the
two endpoints of any cylinder are trivial cofibrations. If \(B\)
is fibrant, both evaluations of any path object are trivial
fibrations.
\end{proposition}
\begin{proof}
Apply the two factorizations of \ref{res:Q01}. For cofibrant \(A\)
each coproduct inclusion into \(A\amalg A\) is a cofibration by
pushout from \(0\to A\). Therefore \(j_i\) is a cofibration,
and \(rj_i=\mathrm{id}\) gives \(j_i\in W\). For fibrant \(B\),
each product projection out of \(B\times B\) is a fibration
by pullback from \(B\to1\). Now \(e_i s=\mathrm{id}\) gives
\(e_i\in W\).
\end{proof}

\begin{theorem}[The homotopy relation]\label{res:Q03}
For cofibrant \(A\) and fibrant \(B\), left homotopy and right
homotopy are the same equivalence relation on maps \(A\to B\).
Left homotopy from \(f\) to \(g\) is a map \(C\to B\) extending
\((f,g)\) on the two ends of a cylinder. Right homotopy is a
map \(A\to P\) with evaluations \(f,g\). Any fixed cylinder
or path object may be used. Homotopy is preserved by
precomposition from a cofibrant object and postcomposition
into a fibrant object. If \(f\) and \(g\) are homotopic,
then \(f\in W\) if and only if \(g\in W\).
\end{theorem}
\begin{proof}
Suppose a left homotopy \(H:C\to B\) is provided. Factor \(r\)
as a trivial cofibration \(C\to C'\) followed by a fibration
\(C'\to A\), which is trivial by two-out-of-three. The first map
allows \(H\) to extend, since \(B\) is fibrant. This gives a
homotopy on a good cylinder. A map from any fixed cylinder to
that good cylinder, preserving both ends and commuting with the
collapse to \(A\), is obtained by lifting. Thus every fixed
cylinder realizes the same left relation. Dually every fixed
path object realizes the same right relation if \(A\) is
cofibrant.

For the equivalence of the relations, choose a good cylinder
and a good path object. A left homotopy \(H\) gives a lift
\(C\to P\) against \(P\to B\times B\) with top map \(s f\)
on the \(j_0\) end and bottom map \((fr,H)\).
Its restriction to \(j_1\) is a right homotopy. Conversely
if \(K:A\to P\) is a right homotopy, lift
\((s f,K):A\amalg A\to P\) against \(e_0:P\to B\)
with bottom map \(fr:C\to B\). Composing the lift with
\(e_1\) gives the left homotopy.

Reflexivity uses \(fr\) and symmetry interchanges the ends.
For transitivity glue two cylinders along the second end
of the first and first end of the second. The outer end
inclusion is a cofibration. It factors through the first
cylinder together with one extra copy of \(A\), after which
one pushes out the boundary inclusion of the second cylinder.
The inclusion of the first cylinder into the glued object
is a pushout of a trivial endpoint cofibration. Hence the
common collapse of the glued object is in \(W\) by
two-out-of-three. The two homotopies agree on the middle end,
so they glue to a homotopy on this cylinder.

Postcomposition of left homotopies is literal. For
precomposition along \(A'\to A\) with \(A'\) cofibrant,
use a good cylinder on \(A\) and lift from a cylinder on
\(A'\), preserving the two new end maps. Finally the
endpoints are both weak equivalences by \ref{res:Q02}.
Apply two-out-of-three to their composites with \(H\).
\end{proof}

\begin{proposition}[Whitehead criterion]\label{res:Q04}
A map between cofibrant fibrant objects is a weak equivalence
if and only if it admits a homotopy inverse.
\end{proposition}
\begin{proof}
For a trivial cofibration \(i:A\to Z\) between such objects,
lift against \(A\to1\) to obtain a retraction \(u:Z\to A\).
Lift along \(i\) into a path object of \(Z\) with endpoints
\((iu,\mathrm{id}_Z)\), using the constant path on \(i\).
This gives \(iu\simeq\mathrm{id}_Z\). Dually a trivial
fibration \(p:Z\to B\) between such objects has a section
which is a homotopy inverse. Factoring an arbitrary weak
equivalence as a trivial cofibration followed by a fibration
gives a trivial fibration by two-out-of-three, with
intermediate object again cofibrant fibrant.

Conversely factor a map having a homotopy inverse as a
cofibration \(i:A\to Z\) followed by a trivial
fibration \(p:Z\to B\). The intermediate object is again
cofibrant fibrant. As \(p\) has a homotopy inverse, so does
\(i\), say \(u_0\). Choose a right homotopy from
\(u_0 i\) to \(\mathrm{id}_A\). Lift it along \(i\) against
the first evaluation from a path object of \(A\), with
bottom map \(u_0\). Its other evaluation gives
\(u:Z\to A\) with \(ui=\mathrm{id}_A\) and \(u\simeq u_0\).
Thus \(iu\simeq\mathrm{id}_Z\), and \(iu\in W\) by
\ref{res:Q03}. The map \(i\) is a retract of the arrow \(iu\)
via the pairs \((i,\mathrm{id}_Z)\) and
\((u,\mathrm{id}_Z)\). Therefore \(i\), and hence \(pi\),
belongs to \(W\).
\end{proof}

\begin{theorem}[Localization by homotopy classes]\label{res:Q05}
The category of cofibrant fibrant objects with morphisms the
homotopy classes is a localization of the model category at
\(W\). A map of the model category becomes an isomorphism
in this localization exactly when it belongs to \(W\).
\end{theorem}
\begin{proof}
Composition is well-defined by \ref{res:Q03}. Choose a
functorial cofibrant replacement \(QX\to X\), a trivial
fibration, and a functorial fibrant replacement \(X\to RX\),
a trivial cofibration. The object \(RQX\) is cofibrant
fibrant. Naturality and two-out-of-three show that the
replacement functors preserve \(W\) and reflect it on
the maps in the two natural replacement squares.
The functor \(X\mapsto RQX\) followed by passage to homotopy
classes inverts \(W\) by \ref{res:Q04}.

For a functor \(F\) inverting \(W\), its restriction to the
cofibrant fibrant objects passes through homotopy classes.
Indeed a cylinder collapse goes to an isomorphism and its two
sections consequently go to the same map. The resulting
factorization of \(F\) is natural up to the natural
isomorphisms induced by \(QX\to X\) and \(QX\to RQX\).
For \(X\) already cofibrant fibrant, \(QX\) is also fibrant
because \(QX\to X\) is a fibration. The two replacement maps
therefore identify, naturally in the homotopy category,
the original cofibrant fibrant object and its new replacement.
Restriction and replacement are thus inverse up to natural
isomorphism on functors that invert \(W\), including natural
transformations. This proves the localization property.
If the image of \(f\) is invertible, the representative
\(RQf\) is a homotopy equivalence, so is in \(W\) by
\ref{res:Q04}. Two-out-of-three in the replacement squares
gives \(f\in W\).
\end{proof}

\begin{proposition}[Cofibrant source and fibrant target]\label{res:Q06}
For cofibrant \(A\) and fibrant \(B\), homotopy classes of
actual maps \(A\to B\) map bijectively to morphisms
\(\operatorname{Ho}(\mathcal C)(A,B)\).
\end{proposition}
\begin{proof}
Replace \(A\) by \(A\hookrightarrow RA\), a trivial
cofibration into a cofibrant fibrant object. Every map to
\(B\) extends. If the restrictions of two extensions are
right homotopic, extend the homotopy along \(A\to RA\)
by lifting against \(P\to B\times B\), with the two
extensions as endpoints. Thus this replacement gives a
bijection on homotopy classes. Replace \(B\) by a trivial
fibration \(QB\to B\), so that \(QB\) is cofibrant
fibrant. For source \(RA\), maps lift against this
trivial fibration. If two composites are left homotopic,
lift the cylinder homotopy with the two maps to \(QB\)
as endpoints. Therefore this replacement also gives a
bijection on classes. Both replacements induce
isomorphisms in the localization of \ref{res:Q05}.
This identifies the stated map with a composite of
bijections. In particular the identification is natural
on maps of cofibrant sources and fibrant targets.
\end{proof}

\section{Quillen adjunctions}

\begin{lemma}[Brown's lemma on cofibrant objects]\label{res:Q07}
A functor on the cofibrant objects which carries trivial
cofibrations to weak equivalences carries every weak
equivalence there to a weak equivalence, assuming two-out-of-three
in the target. The corresponding assertion for fibrant
objects and trivial fibrations holds dually.
\end{lemma}
\begin{proof}
Let \(f:X\to Y\) be a weak equivalence of cofibrant objects.
Factor \((f,\mathrm{id}_Y):X\amalg Y\to Y\) as a
cofibration followed by a trivial fibration \(p:Z\to Y\).
Both maps \(X\to Z\) and \(Y\to Z\) are trivial
cofibrations by two-out-of-three, and \(Z\) is cofibrant.
The image of \(p\) is weak because its composite with
the image of \(Y\to Z\) is the identity. The same then
holds for the image of \(f\).
\end{proof}

An adjunction \(L:\mathcal C\rightleftarrows\mathcal D:U\)
is Quillen when \(L\) preserves cofibrations and trivial
cofibrations, equivalently when \(U\) preserves fibrations
and trivial fibrations. This equivalence follows by
transporting lifting squares along the adjunction and
using the orthogonality observation after \ref{res:Q01}.

\begin{theorem}[Derived Quillen adjunction]\label{res:Q08}
A Quillen adjunction induces an adjunction
\[
 \mathbf L L:\operatorname{Ho}(\mathcal C)
 \rightleftarrows \operatorname{Ho}(\mathcal D):\mathbf R U.
\]
It is computed by \(LQ\) and \(UR\), where \(Q\) is
cofibrant replacement in \(\mathcal C\) and \(R\)
is fibrant replacement in \(\mathcal D\).
\end{theorem}
\begin{proof}
The two composites preserve weak equivalences by
\ref{res:Q07} and hence descend by \ref{res:Q05}.
The functors \(L\) and \(U\) preserve respectively
cofibrancy and fibrancy because they preserve the relevant
structure maps and initial or terminal objects.
For cofibrant \(X\) and fibrant \(Y\), the ordinary adjunction
identifies maps \(LX\to Y\) with maps \(X\to UY\).
Apply \(L\) to a cylinder of \(X\). Its end inclusion
is a cofibration and the collapse is a weak equivalence
by \ref{res:Q07}, since both ends of the collapse map
are cofibrant. It is a cylinder of \(LX\) because \(L\)
preserves the coproduct. The ordinary adjunction identifies
maps out of this cylinder with maps out of the original
cylinder on the other side and respects endpoint maps.
By \ref{res:Q03}, the adjunction therefore identifies
the two homotopy relations. Applying \ref{res:Q06} gives
the adjunction on homotopy categories. Replacing both
variables supplies this bijection for arbitrary objects,
naturally by the functorial choices and the natural
identification of \ref{res:Q06}.
\end{proof}

\begin{theorem}[Quillen equivalence criterion]\label{res:Q09}
The derived adjunction of \ref{res:Q08} is an equivalence
if and only if the following condition holds. For every
cofibrant \(X\) and fibrant \(Y\), a map \(LX\to Y\) is
a weak equivalence exactly when its adjoint \(X\to UY\)
is a weak equivalence.
\end{theorem}
\begin{proof}
Assume the condition. For cofibrant \(X\), choose a
fibrant replacement \(LX\to R LX\). Its adjoint is
weak by the condition, and represents the derived unit.
For fibrant \(Y\), choose a cofibrant replacement
\(Q UY\to UY\). Its adjoint is weak and represents
the derived counit. Thus the unit and counit are
isomorphisms by \ref{res:Q05}, and functorial
replacement gives this conclusion on all objects.
Conversely, an adjunction which is an equivalence
identifies isomorphisms with isomorphisms under the
adjunction bijection. Apply \ref{res:Q05} to both
maps in the claimed criterion. This proof uses the
unit at a fibrant replacement of \(LX\) and the
counit at a cofibrant replacement of \(UY\), in agreement
with \cite[correction to Corollary 1.3.16]{hovey-errata}.
\end{proof}


\section{Cochain complexes over a ring}

Let \(A\) be a ring with identity. Write \(\operatorname{Ch}(A)\)
for the category of unbounded cochain complexes of left \(A\)-modules.
It has all small limits and colimits, obtained degree by degree.
We use quasi-isomorphisms as weak equivalences and reindex the
projective disk generators of \cite[Section 2.3]{hovey} cohomologically.

\begin{proposition}[Spheres, disks, and their lifts]\label{res:Q10}
Let \(S^n\) be \(A\) in degree \(n\) and \(D^n\) be
\(A\xrightarrow{\mathrm{id}} A\) in degrees \(n,n+1\).
Use the sets of maps
\[
 I=\{S^{n+1}\longrightarrow D^n:n\in\mathbb Z\},
 \qquad
 J=\{0\longrightarrow D^n:n\in\mathbb Z\},
\]
where \(S^{n+1}\to D^n\) includes the upper degree. A map has
the right lifting property against \(J\) precisely when it is
degreewise surjective. It has the right lifting property against
\(I\) precisely when it is a degreewise surjective quasi-isomorphism.
\end{proposition}
\begin{proof}
A map out of \(D^n\) specifies any element in degree \(n\).
Lifting against \(J\) thus tests degreewise surjectivity.
A square against \(S^{n+1}\to D^n\) for \(p:X\to Y\) specifies
\(z\in Z^{n+1}(X)\) and \(y\in Y^n\) such that \(p(z)=dy\).
A lift is \(x\in X^n\) with \(p(x)=y\) and \(dx=z\).
Taking \(z=0\) gives surjectivity on cycles in degree \(n\).
For general \(y\), first lift the cycle \(dy\) and then
lift \(y\), obtaining degreewise surjectivity. Taking \(y=0\)
shows that the kernel has zero cohomology.

Conversely suppose \(p\) is degreewise surjective and a
quasi-isomorphism. Its kernel is acyclic by the long exact
sequence of cohomology. Start with a graded lift \(x_0\) of
\(y\). The difference \(z-dx_0\) is a cycle in the kernel,
hence the differential of some \(u\) in the kernel. Then
\(x=x_0+u\) gives the required lift.
\end{proof}

\begin{lemma}[Free cell factorizations]\label{res:Q11}
Every cochain map factors functorially both as an \(I\)-cell
map followed by an \(I\)-right lifting map and as a \(J\)-cell
map followed by a \(J\)-right lifting map. Here a cell is
attached by pushing out a coproduct of generators, and a
cell map is a sequential composite of such attachments.
An \(I\)-cell map is obtained by adjoining free graded
generators with differential in the previous stage.
A \(J\)-cell map adjoins a direct sum of disks.
\end{lemma}
\begin{proof}
Start with the given arrow \(X_0\to Y\). Given \(X_t\to Y\),
take the coproduct of one copy of each generator for every
commutative square from that generator to \(X_t\to Y\).
Push out its domain map to obtain \(X_{t+1}\), with the
induced map to \(Y\). In degree \(n\), an \(I\)-attachment
adjoins one free generator \(e\) of degree \(n\) for each
corresponding square, imposing \(de=z\) where \(z\) is the
image of the sphere generator. This is consistent because
\(z\) is closed. For \(J\), the whole disk is added freely
as a direct summand of complexes. Take the sequential
colimit \(X_\infty\). The domain of an \(I\)-generator is
one cycle generator. A map from it to \(X_\infty\) is
represented at a finite stage because the cycle condition
becomes true at some finite stage. A \(J\)-domain is zero.
Thus any square at the colimit factors through a finite
stage, after which it acquires its filler. Using all
squares at each stage makes the factorizations functorial.
\end{proof}

\begin{theorem}[The unbounded projective model structure]\label{res:Q12}
There is a model structure on \(\operatorname{Ch}(A)\)
with weak equivalences the quasi-isomorphisms, fibrations
the degreewise surjections, and cofibrations the maps
lifting against degreewise surjective quasi-isomorphisms.
Every object is fibrant. Its homotopy category is
\(\mathrm D(A)\).
\end{theorem}
\begin{proof}
The weak equivalences have two-out-of-three and are
retract-closed by functorial cohomology. The lifting-defined
classes are retract-closed. By \ref{res:Q10}, the two
factorizations of \ref{res:Q11} have the required right
classes. Every \(I\)-cell map is a cofibration by its
lifting property. Every \(J\)-generator is an \(I\)-cell
map, since \(0\to S^{n+1}\) is the pushout of
\(S^{n+2}\to D^{n+1}\) along the zero map and
\(S^{n+1}\to D^n\) is itself a generator. Consequently
\(J\)-cell maps are cofibrations. They are also
quasi-isomorphisms since adjoining a direct sum of disks
does not change cohomology.

It remains to identify the left lifting class against
fibrations. Every map in this class is a retract of its
\(J\)-cell factor by the lift against the right factor.
It is therefore a trivial cofibration. Conversely
factor a trivial cofibration \(i\) as \(i=pj\) with \(j\)
a \(J\)-cell map and \(p\) a fibration. As \(i\) and \(j\)
are quasi-isomorphisms, \(p\) is a trivial fibration.
The lift furnished by \(i\) exhibits \(i\) as a retract
of \(j\), which has the left lifting property against
every fibration. This verifies both lifting classes and
all the axioms of \ref{res:Q01}. All maps to the zero complex
are surjective. The final assertion follows from \ref{res:Q05}
and the definition of the derived category by localization.
\end{proof}

\begin{proposition}[Cofibrations and acyclic targets]\label{res:Q13}
A map of complexes is a cofibration in \ref{res:Q12}
if and only if it is a degreewise split injection whose
cokernel \(P\) is degreewise projective and \(K\)-projective.
Here \(K\)-projective means that
\(\operatorname{Hom}_A^\bullet(P,E)\) is acyclic for every
acyclic \(E\). The cokernel is also \(K\)-flat, meaning
that \(E\otimes_A P\) is acyclic for every acyclic complex
of right \(A\)-modules.
\end{proposition}
\begin{proof}
A cofibration is a retract of its \(I\)-cell factor.
An \(I\)-cell map is split injective as a graded map,
and its cokernel has a filtration with graded-split
inclusions and freely adjoined generators whose
differentials land in earlier stages. A cochain map \(f\)
from such a cokernel to an acyclic \(E\) is null homotopic.
At the stage adjoining \(e\), with \(de\) already present,
the equation \(f(e)=d h(e)+h(de)\) is solvable because
\(f(e)-h(de)\) is a cycle. Indeed
\(d f(e)-d h(de)=f(de)-f(de)+h(d^2e)=0\)
using the equation on the prior stage. Choose these values
on all generators and extend linearly. The same argument
applied to all shifts of \(E\) is precisely acyclicity
of the Hom complex. The cokernel is \(K\)-flat because
tensoring the graded-split filtration with \(E\) gives
a filtration by complexes with acyclic quotients, direct
sums of shifts of \(E\). Each finite stage is acyclic
by the long exact sequence and the union is acyclic
by exactness of filtered colimits of modules. Taking
retracts preserves all these conditions, which proves
necessity.

For sufficiency let \(i:A_0\to B_0\) have the stated
properties and consider a lifting square against a
surjective quasi-isomorphism \(q:X\to Y\). Using the
graded splitting of \(i\) and degreewise projectivity
of \(P=B_0/A_0\), choose a graded lift \(u:B_0\to X\)
extending the prescribed map on \(A_0\).
Its defect \(d_Xu-u d_{B_0}\) vanishes on \(A_0\)
and has image in \(\ker q\). It defines a degree-one
cocycle of \(\operatorname{Hom}_A^\bullet(P,\ker q)\).
The identity on two successive differentials gives the
cocycle equation. By \ref{res:Q10}, the kernel is
acyclic. Thus the defect is the Hom differential of
a degree-zero map \(P\to\ker q\). Subtract this
map after quotient and inclusion to correct \(u\)
to a cochain map and complete the lift.
\end{proof}

\begin{corollary}[Chain homotopy in the projective model]\label{res:Q14}
For cofibrant \(X\), the morphisms from \(X\) to any
complex \(Y\) in \(\mathrm D(A)\) are cochain maps modulo
ordinary cochain homotopy.
\end{corollary}
\begin{proof}
Use the path complex
\[
 P(Y)^n=Y^n\oplus Y^n\oplus Y^{n-1},\qquad
 d(y_0,y_1,b)=(dy_0,dy_1,y_1-y_0-db).
\]
Its third component after two applications of \(d\) is
zero, and its evaluations are the two projections. They
together form a degreewise surjection to \(Y\times Y\).
The first evaluation has contractible kernel, since in
coordinates \((a,b)\) that kernel has differential
\((da,a-db)\) and contracting homotopy \((a,b)\mapsto(b,0)\).
The diagonal \(y\mapsto(y,y,0)\) is therefore a
quasi-isomorphism. A cochain map \(X\to P(Y)\) with
first two coordinates \(f,g\) has third coordinate
\(h\) with \(g-f=d h+h d\). Now apply \ref{res:Q03}
and \ref{res:Q06}, using fibrancy of all complexes.
\end{proof}

\section{Tensor products and homotopy limits}

\begin{theorem}[Tensor compatibility]\label{res:Q15}
If \(A\) is commutative, the projective model structure
is a closed symmetric monoidal model structure for the
usual tensor product of cochain complexes. In particular,
the pushout product of two cofibrations is a cofibration
and is trivial if either is trivial.
\end{theorem}
\begin{proof}
The tensor product uses direct sums in total degree and
the signed differential. Its internal Hom uses products
in each degree with differential
\(d(h)=d h-(-1)^{|h|}h d\). Evaluation and currying give
the adjunction by expansion of the cochain identities.
Symmetry is \(x\otimes y\mapsto(-1)^{|x||y|}y\otimes x\).

For \(a:A_0\to B_0\), \(b:C_0\to D_0\), their
pushout product is
\[
 (B_0\otimes C_0)\amalg_{A_0\otimes C_0}(A_0\otimes D_0)
 \longrightarrow B_0\otimes D_0.
\]
For \(a,b\in I\), the inclusion has the product of the
two lower disk generators as its only new generator.
Its differential lies in the source, so this is a
single \(I\)-attachment. For \(a\in I\), \(b\in J\),
the inclusion splits degreewise and its quotient is
a shift of a disk, from the lower sphere of the
\(I\)-cokernel tensored with the disk of \(b\).
The sign is absorbed by changing the sign of one disk
generator. The quotient is acyclic, degreewise free
and \(K\)-projective by contractibility. Thus the
pushout product is a trivial cofibration by
\ref{res:Q13}. Symmetry handles \(I,J\) in the other order.

For fixed \(a\), the operation \(a\square -\)
preserves pushouts of arrows by separate cocontinuity
of the tensor product. For a composite \(C\to D\to E\)
in the other variable, first pass through
\[
 (B_0\otimes D)\amalg_{A_0\otimes D}(A_0\otimes E)
\]
on the way to \(B_0\otimes E\). The first map in
this factorization is a pushout of the pushout product
for \(C\to D\), and the second is the pushout product
for \(D\to E\). For a sequence of attachments use these
same intermediates with its final object \(E\), taking
colimits at limit stages. This proves closure under
the cell operations in either variable. Retracts are
preserved because the pushout product is functorial
on pairs of arrows. Since cofibrations are retracts
of \(I\)-cell maps and trivial cofibrations retracts
of \(J\)-cell maps by \ref{res:Q12}, the generator
calculations prove the assertion. The unit \(A[0]\)
is cofibrant by \ref{res:Q13}, so the identity unit
replacement verifies the unit axiom. For any weak
cofibrant replacement of this unit the same conclusion
follows from \ref{res:Q07} applied to tensor with a
cofibrant object, which is left Quillen by the
pushout-product statement.
\end{proof}

Use the cone convention
\(\operatorname{Cone}(h)^n=Z^n\oplus T^{n+1}\)
for \(h:T\to Z\), with differential
\(d(z,t)=(dz+h(t),-dt)\).
The cochain shift \([-1]\) negates the differential.
For the general homotopy-limit viewpoint see
\cite[Chapters 5--6]{riehl}. The explicit models below use exact
products in module categories, as in \cite[Section 34]{st-derived}.

\begin{proposition}[Homotopy pullbacks of complexes]\label{res:Q16}
For a span \(X\xrightarrow{f} Z\xleftarrow{g}Y\),
a homotopy pullback is
\[
 \operatorname{Cone}(f-g:X\oplus Y\longrightarrow Z)[-1].
\]
It is functorial on strict maps of spans and invariant under
objectwise quasi-isomorphism. When either leg is a fibration,
the strict pullback maps to this model by a quasi-isomorphism.
\end{proposition}
\begin{proof}
Put \(C^n=Z^n\oplus Z^{n-1}\) with differential
\((a,b)\mapsto(da,a-db)\). The homotopy
\((a,b)\mapsto(b,0)\) contracts \(C\). Replace \(f\)
by the surjection \(X\oplus C\to Z\) given by
\((x,a,b)\mapsto f(x)+a\). The original \(X\) includes
by a quasi-isomorphism over \(Z\). The strict pullback
of the resulting span consists of triples
\((x,y,b)\) with \(a=g(y)-f(x)\), and its differential
is \((dx,dy,g(y)-f(x)-db)\). This is precisely the
stated cone with the coordinates reordered.
A strict objectwise quasi-isomorphism of spans induces
one on cones by the cone cohomology sequence.
If a leg was already surjective, then \(f-g\)
is surjective and its kernel is the ordinary pullback.
The short exact sequence with this kernel identifies
its inclusion into the shifted cone as a quasi-isomorphism.
These statements show that fibrational replacement
of one leg followed by ordinary pullback gives a
homotopy-invariant model and that this agrees with
ordinary pullback on spans having that leg fibrant.
\end{proof}

\begin{theorem}[Sequential homotopy limits]\label{res:Q17}
For a tower \(X_0\leftarrow X_1\leftarrow\cdots\) of
complexes, let \(s:\prod_n X_n\to\prod_n X_n\) carry the
\((n+1)\)-st coordinate to the \(n\)-th by the transition map.
A model for \(\operatorname{holim}_n X_n\) is
\[
 \operatorname{Cone}\bigl(1-s:\prod_n X_n\to\prod_n X_n\bigr)[-1].
\]
If all transition maps are degreewise surjective, the
ordinary inverse limit maps quasi-isomorphically to this
model. In each degree \(q\) there is a natural short
exact sequence
\[
 0\longrightarrow\lim\nolimits^1_n H^{q-1}(X_n)
 \longrightarrow H^q(\operatorname{holim}_n X_n)
 \longrightarrow\lim_n H^q(X_n)\longrightarrow0 ,
\]
where for a sequential system of modules
\(\lim^1\) denotes the cokernel of \(1-s\) on their product.
\end{theorem}
\begin{proof}
Products of modules are exact. Taking the cone therefore
gives an objectwise quasi-isomorphism invariant construction.
For a degreewise surjective tower, given a tuple \((b_n)\),
choose \(a_0=0\) and solve
\(s_n(a_{n+1})=a_n-b_n\) recursively in each degree.
Hence \(1-s\) is degreewise surjective and its kernel
is the ordinary inverse limit. The short exact sequence
of complexes proves the asserted comparison with the shifted
cone. Every tower admits a functorial levelwise
quasi-isomorphism to such a tower. Start with
\(Y_0=X_0\). At each subsequent stage factor the composite
\(X_n\to X_{n-1}\to Y_{n-1}\) as a trivial cofibration
\(X_n\to Y_n\) followed by a fibration \(Y_n\to Y_{n-1}\)
using \ref{res:Q12}. Thus taking the ordinary limit
after this replacement is equivalent to the cone model.
Taking cohomology of the cone exact sequence gives the
displayed short exact sequence, since cohomology commutes
with products and the kernel of \(1-s\) on a product
is the ordinary inverse limit.
\end{proof}

\paragraph{Scope.}
The construction above concerns the projective model structure on
module complexes. The unbounded injective resolutions used for
derived Hom are constructed in the derived-functor chapter. The
existence of an injective model structure on sheaf complexes has
different hypotheses and does not enter these proofs.
